\documentclass[11pt]{article}
\usepackage[letterpaper,margin=1in]{geometry}
\usepackage[T1]{fontenc}
\usepackage{lmodern,microtype,amsmath,amssymb,amsthm}
\usepackage{xcolor}
\usepackage[colorlinks=true,linkcolor=black,citecolor=blue!45!black,urlcolor=blue!45!black]{hyperref}
\usepackage{xurl}
\newtheorem{theorem}{Theorem}[section]
\newtheorem{lemma}[theorem]{Lemma}
\newtheorem{corollary}[theorem]{Corollary}
\theoremstyle{definition}
\newtheorem{definition}[theorem]{Definition}
\theoremstyle{remark}
\newtheorem{remark}[theorem]{Remark}
\newcommand{\F}{\mathbb{F}}
\newcommand{\Z}{\mathbb{Z}}
\newcommand{\Tr}{\operatorname{Tr}}
\newcommand{\Cx}{C_{\oplus}}
\newcommand{\supp}{\operatorname{supp}}
\newcommand{\rk}{\operatorname{rank}}
\hypersetup{pdftitle={Superlinear XOR circuit lower bounds and superlinear rigidity for an explicit matrix over F2},
  pdfauthor={Charlie Cruz},
  pdfkeywords={superlinear lower bounds, XOR circuits, linear circuits, matrix rigidity, F2, GF(2), explicit matrix, Fourier transform, circuit complexity, trace descent}}
\title{Superlinear XOR Circuit Lower Bounds and\\Superlinear Rigidity for an Explicit Matrix over $\F_2$}
\author{Charlie Cruz}
\date{November 24, 2026}
\begin{document}
\maketitle
\begin{abstract}
An \emph{XOR circuit} computes a linear map $x\mapsto Ax$ over $\F_2$ using
only fan-in-two XOR gates. The \emph{rigidity} of $A$ at rank $r$ is the
least number of entries of $A$ that must be changed to bring its rank down
to $r$. We give an explicit $0/1$ matrix $H_N\in\F_2^{N\times N}$ for which
both quantities are superlinear. Every XOR circuit computing $H_N$ has
$\Omega(N\sqrt{\log N})$ gates, while $N(\log N)^{1/2+o(1)}$ gates suffice.
Lowering the rank of $H_N$ over $\F_2$ to $\gamma N$ requires changing more
than $N2^{\Omega(\sqrt{\log N})}$ entries. Both bounds hold simultaneously, so
$H_N$ is not a small XOR circuit plus a low-rank matrix plus a sparse matrix.
They also hold for direct sums of many copies and for XOR circuits with
feedback. The matrix $H_N$ is the Fourier transform of $(\Z/3\Z)^d$ over
$\F_4$, written in binary. It is equally hard over every field of
characteristic two, and the Walsh--Hadamard matrix plays the same role
over every other field. The proofs reduce to the recent theorem that
Fourier transforms over fields with enough roots of unity need
$\Omega(N\sqrt{\log N})$ linear gates. That theorem never applies over
$\F_2$, and its proof has been formally verified in Lean. Two elementary
tools carry it down to $\F_2$. One is a trace descent for circuits. The
other is a projection onto $\F_4$-linear maps, which controls rank and
sparsity. To our knowledge these are the first superlinear lower bounds
for an explicit $\F_2$-linear map against XOR circuits of unbounded depth.
The rigidity bound is also the first explicit bound over $\F_2$ that is
superlinear at linear rank.
\end{abstract}

\section{Introduction}

\subsection{Background}
An \emph{XOR circuit} computes a linear map $\F_2^N\to\F_2^N$. Each gate
adds two previously computed bits, and each output is an input or a gate.
Write $\Cx(A)$ for the least number of gates computing $x\mapsto Ax$.
Lupanov showed that $\Cx(A)=O(N^2/\log N)$ for every $A$ \cite{Lup}, and
counting shows this is tight for almost all $A$. For explicit matrices,
however, no superlinear lower bound against XOR circuits of unbounded
depth has been known. Superlinear bounds are known only for restricted
models, such as bounded depth.

Valiant proposed matrix rigidity as a route to such bounds for circuits
of logarithmic depth \cite{Val77}. The rigidity $R_A(r)$ is the least
number of entries of $A$ that must be changed to bring its rank down to at
most $r$. Valiant showed that $R_A(\varepsilon N)\ge N^{1+\delta}$ implies
superlinear lower bounds for log-depth linear circuits computing $A$.
Over finite fields, the best rigidity lower bound for an explicit matrix
is $\Omega((N^2/r)\log(N/r))$ \cite{Fri93,SSS97}, which is only linear at
rank $r=\Theta(N)$. Natural candidates also fall short of Valiant's
threshold. For example, Hadamard matrices are not rigid enough
\cite{AW17}.

\subsection{The Fourier exponent theorem}
Recently, linear circuit lower bounds for an explicit family were proved
over a large class of fields \cite{FFT}. For a field $K$, a \emph{linear
circuit} over $K$ has gates computing $au+bv$, with $a,b\in K$ arbitrary
and $u,v$ previously computed. An output may select a single input or gate
and multiply it by a scalar. A \emph{feedback system} with $s$ states is a
system of equations $v=Cv+Ax$ in which each state is a weighted sum of at
most two inputs or states. The matrix $I-C$ must be invertible, and the
outputs are read as in a circuit. Every circuit with $s$ gates is such a
system. Write $C_K(T)$ for the least number of gates in a linear circuit
over $K$ computing $T$, and $F_G$ for the character table (Fourier
transform) of a finite abelian group $G$.

\begin{theorem}[{\cite[Theorems~1.1 and~1.2]{FFT}}]\label{thm:fft}
There are absolute constants $c,\beta,\eta>0$ such that the following
holds. Let $G$ be a finite abelian group of order $n\ge2$, and let $K$ be
a field in which $n$ is invertible and which contains a primitive root of
unity of order $\exp(G)$. Let $T=F_G^{\oplus J}$ for some $J\ge1$, and put
$h=Jn$.
\begin{enumerate}
\item $c\,h\sqrt{\log n}\le C_K(T)\le O_\varepsilon\bigl(h(1+\log n)^{1/2+\varepsilon}\bigr)$
for every $\varepsilon>0$.
\item Suppose $T=M+B+Z$ over $K$. Here $M$ is computed by a feedback
system with $s$ states, $\rk B\le\beta h$, and
$|\supp Z|\le\eta h\,2^{\sqrt{\log n}/8}$. Then $s\ge c\,h\sqrt{\log n}$.
\end{enumerate}
\end{theorem}

In \cite{FFT}, part~2 is proved even for $T$ in the Zariski closure of
such sums. Taking $M=0$ in part~2 gives a rigidity bound for $F_G$. The
hypotheses never hold over $\F_2$, because the only root of unity there is
$1$. The natural candidate, the Walsh--Hadamard matrix, has entries
$\pm1$, which all reduce to $1$ modulo two.

\subsection{Results}
Let $\omega\in\F_4$ be a root of $\omega^2=\omega+1$, and identify $\F_4$
with $\F_2^2$ through the basis $\{1,\omega\}$. Multiplication by $\omega$
then acts by
\[
W=\begin{pmatrix}0&1\\1&1\end{pmatrix},\qquad W^3=I.
\]

\begin{definition}\label{def:H}
For $d\ge1$, let $n=3^d$ and $N=2n$. Let $H_N\in\F_2^{N\times N}$ be the
$n\times n$ block matrix with $2\times2$ blocks
\[
(H_N)_{a,g}=W^{\langle a,g\rangle\bmod3},\qquad a,g\in(\Z/3\Z)^d.
\]
For other $N\ge6$, let $d$ be maximal with $2\cdot3^d\le N$, and define
$H_N=H_{2\cdot3^d}\oplus I_{N-2\cdot3^d}$.
\end{definition}

For example,
\[
H_6=\begin{pmatrix}I&I&I\\I&W&W^2\\I&W^2&W\end{pmatrix}.
\]
For $N=2\cdot3^d$, the matrix $H_N$ is the Fourier transform of
$(\Z/3\Z)^d$ over $\F_4$, written over $\F_2$. Its entries can be computed
in polynomial time.

Our main result is a robust lower bound over $\F_2$.

\begin{theorem}\label{thm:main}
There are absolute constants $c',\gamma,\eta'>0$ such that the following
holds for all sufficiently large $N=2\cdot3^d$ and all $J\ge1$. Suppose
that over $\F_2$,
\[
H_N^{\oplus J}=M+B+Z,
\]
where $M$ is computed by an XOR feedback system with $s$ states,
$\rk B\le\gamma JN$, and $|\supp Z|\le\eta'JN\,2^{\sqrt{\log N}/8}$. Then
$s\ge c'JN\sqrt{\log N}$.
\end{theorem}

The two extreme cases are worth stating on their own.

\begin{corollary}[XOR circuits]\label{cor:xor}
For all $N\ge6$, we have $\Cx(H_N)=N(\log N)^{1/2+o(1)}$, and
$\Cx(H_N)\ge c'N\sqrt{\log N}$ for all sufficiently large $N$. Moreover,
$\Cx(H_N^{\oplus J})\ge c'JN\sqrt{\log N}$, even though the circuit may mix
the copies.
\end{corollary}

\begin{corollary}[Rigidity over $\F_2$]\label{cor:rigid}
For all sufficiently large $N=2\cdot3^d$,
\[
R_{H_N}(\gamma N)>\eta'N\,2^{\sqrt{\log N}/8}.
\]
\end{corollary}

The rigidity bound is superlinear at linear rank. It is below Valiant's
threshold $N^{1+\delta}$, since $2^{\sqrt{\log N}}=N^{o(1)}$, so it does not
yield lower bounds for log-depth circuits.

The same tools give explicit hard matrices over every field.

\begin{theorem}\label{thm:every}
Let $K$ be a field.
\begin{enumerate}
\item If $\operatorname{char}K\ne2$, the Walsh--Hadamard matrix of order
$N=2^d$ satisfies $C_K=N(\log N)^{1/2+o(1)}$ and $C_K\ge c''N\sqrt{\log N}$.
\item If $\operatorname{char}K=2$, the $0/1$ matrix $H_N$ satisfies
$C_K(H_N)=N(\log N)^{1/2+o(1)}$ and $C_K(H_N)\ge c''N\sqrt{\log N}$.
\end{enumerate}
Here $c''>0$ is absolute.
\end{theorem}

Part~1 is Theorem~\ref{thm:fft} with $G=(\Z/2\Z)^d$, since $2^d$ is
invertible and $-1$ is a primitive square root of unity. Part~2 shows that
a single $0/1$ matrix is hard uniformly over all fields of characteristic
two.

\subsection{Techniques}
Two elementary observations carry Theorem~\ref{thm:fft} from $\F_4$ down
to $\F_2$.

\emph{Descent for circuits} (Lemma~\ref{lem:descent}). Every
$\F_2$-linear functional on $\F_4^n$ is the trace of a unique
$\F_4$-linear form, and this correspondence is $\F_2$-linear. Replacing
every wire of an XOR circuit by the corresponding form therefore turns an
XOR circuit for the binary form of an $\F_4$-linear map into an
$\F_4$-linear circuit for the map. The only cost is one scaling gate per
input bit.

\emph{Projection for matrices} (Lemma~\ref{lem:proj}). In characteristic
two, $3$ is invertible. So averaging over conjugation by multiplication by
$\omega$ projects every $\F_2$-linear map onto its $\F_4$-linear part. The
projection fixes $H_N$ and at most triples rank. It never spreads a
nonzero entry outside its own $2\times2$ block, and it maps small XOR
systems to small XOR systems. Applying it to a decomposition
$H_N=M+B+Z$ gives a decomposition of the Fourier transform over $\F_4$
with comparable budgets. Theorem~\ref{thm:fft} forbids such a
decomposition.

\subsection{Status of the imported theorem}
Everything beyond Theorem~\ref{thm:fft} is proved in full below. The full
statement of Theorem~\ref{thm:fft} appears as the Lean theorem
\texttt{batchedFourierComplexity} in the challenge file
\texttt{ExactFourier.lean} of the accompanying formalization \cite{Lean}.
That statement covers batches, feedback systems, and rank and support
budgets $\beta h$ and $\eta h2^{\sqrt{\log_2n}/8}$, and its comparator
configuration certifies a proof in \texttt{OAI.LinearAlgebra.FiniteFourier.Main}
using only the standard axioms \texttt{propext}, \texttt{Quot.sound} and
\texttt{Classical.choice}. We have checked that the Lean statement matches
Theorem~\ref{thm:fft}, but we have not rerun the comparator.

We have also read the paper proof of the lower bound (Sections~2--3 of
\cite{FFT}) line by line, and found it correct given the properties of
planar spectral prices from \cite[Theorem~1.1]{Geom}. In \cite{Geom},
those properties rest on the existence of spectral flats. That proof is
long and we have not verified it, but the formalization constructs the
prices directly \cite{Lean}.

\section{Preliminaries}
All circuits have fan-in two. An XOR circuit is a linear circuit over
$\F_2$, and an XOR feedback system is a feedback system over $\F_2$. For a
matrix $X$, $\supp X$ is the set of its nonzero entries. We write
$\Phi_n=I_n\otimes W$ for multiplication by $\omega$ on
$\F_4^n\cong\F_2^{2n}$. An $\F_2$-linear map $X$ on $\F_2^{2n}$ is
\emph{$\F_4$-linear} if it commutes with $\Phi_n$. We then also regard
$X$ as an $n\times n$ matrix over $\F_4$, whose $(i,j)$ entry corresponds
to the $(i,j)$ block of $X$.

For a finite separable extension $E/K$ of degree $k$, fix a $K$-basis of
$E$. The \emph{$K$-form} $T_K\in K^{kn\times kn}$ of an $E$-linear map
$T:E^n\to E^n$ is $T$ written in coordinates over that basis.

\section{Descent}
\begin{lemma}\label{lem:descent}
Let $E/K$ be a finite separable extension of degree $k$, and let
$T:E^n\to E^n$ be $E$-linear. If $T_K$ is computed by a linear circuit (or
feedback system) over $K$ with $s$ gates (or states), then $T$ is computed
by one over $E$ with at most $s+kn$. Conversely,
$C_K(T_K)\le O_E(C_E(T)+n)$.
\end{lemma}

\begin{proof}
Since $E/K$ is separable, the trace form $(c,x)\mapsto\Tr_{E/K}(cx)$ is
nondegenerate. Hence every $K$-linear functional $\phi:E^n\to K$ has the
form $\phi(x)=\Tr(\sum_ic_ix_i)$ for a unique $c\in E^n$. Let
$\hat\phi(x)=\sum_ic_ix_i$ be its \emph{lift}. Lifting is $K$-linear:
$\widehat{a\phi+b\psi}=a\hat\phi+b\hat\psi$ for $a,b\in K$.

Let $\theta_1^*,\dots,\theta_k^*$ be the basis dual to the fixed basis
under the trace form. Then the input coordinates of $T_K$ are the
functionals $\Tr(\theta_j^*x_i)$, whose lifts are $\theta_j^*x_i$. Given a
circuit for $T_K$, replace every wire $\phi$ by $\hat\phi$:
\begin{itemize}
\item each input becomes $\theta_j^*x_i$, at the cost of one scaling gate,
for $kn$ gates in all;
\item each gate $a\phi+b\psi$ becomes the gate $a\hat\phi+b\hat\psi$;
\item an output computing the coordinate $\Tr(\theta_j^*T(x)_i)$ becomes
$\theta_j^*T(x)_i$, and the readout rescales it to $T(x)_i$ at no cost.
\end{itemize}
For a feedback system $v=Cv+Ax$, the lifted system has the same matrix
$C$ with entries in $K$. Since $I-C$ is invertible over $K$, it is
invertible over $E$, so the lifted system is uniquely solvable and its
solution consists of the lifts.

For the converse, each gate over $E$ is a fixed $K$-linear map on $2k$
coordinates, which costs $O_E(1)$ gates over $K$.
\end{proof}

\section{Projection}
\begin{lemma}\label{lem:proj}
For $X\in\F_2^{2n\times2n}$ put
\[
\Pi(X)=\sum_{k=0}^{2}\Phi_n^{-k}X\,\Phi_n^{k}.
\]
Then:
\begin{enumerate}
\item $\Pi(X)$ is $\F_4$-linear, and $\Pi(X)=X$ whenever $X$ is
$\F_4$-linear;
\item $\rk_{\F_4}\Pi(X)\le\tfrac32\rk_{\F_2}X$;
\item as a matrix over $\F_4$, $\Pi(X)$ has at most $|\supp X|$ nonzero
entries;
\item if $X$ is computed by an XOR feedback system with $s$ states, then
$\Pi(X)$ is computed by one with at most $3s+O(n)$ states.
\end{enumerate}
\end{lemma}

\begin{proof}
(1) Conjugation by $\Phi_n$ permutes the three summands cyclically, so
$\Pi(X)$ commutes with $\Phi_n$. If $X$ commutes with $\Phi_n$, all
three summands equal $X$, and $3X=X$ over $\F_2$.

(2) Each summand has $\F_2$-rank $\rk X$, so
$\rk_{\F_2}\Pi(X)\le3\rk X$. An $\F_4$-linear map of $\F_2$-rank $\rho$ has
$\F_4$-rank $\rho/2$.

(3) Since $\Phi_n$ is block diagonal with $2\times2$ blocks, each summand
is supported on the blocks that meet $\supp X$. There are at most
$|\supp X|$ such blocks, and each nonzero block of $\Pi(X)$ is a single
nonzero entry over $\F_4$.

(4) Multiplying a $2$-bit block by $W$ or by $W^{-1}=W^2$ costs one XOR
gate. Run three copies of the system, the $k$-th with $\Phi_n^k$ applied
at its inputs and $\Phi_n^{-k}$ at its outputs, and add the three
outputs with $2\cdot2n$ further gates.
\end{proof}

\section{Proofs}
\begin{proof}[Proof of Theorem~\ref{thm:main}]
Let $n=3^d$, $h=Jn$, and let $T$ be the direct sum of $J$ copies of the
Fourier transform of $(\Z/3\Z)^d$ over $\F_4$. The order $3^d$ is odd and
$\omega$ is a primitive cube root of unity, so Theorem~\ref{thm:fft}
applies. The $\F_2$-form of $T$ is $H_N^{\oplus J}$, which is $\F_4$-linear.

Suppose $H_N^{\oplus J}=M+B+Z$ as in the statement. By
Lemma~\ref{lem:proj}(1),
\[
T=\Pi(M)+\Pi(B)+\Pi(Z),
\]
with all three terms $\F_4$-linear. The three terms satisfy the following
bounds.
\begin{itemize}
\item By Lemma~\ref{lem:proj}(4) and Lemma~\ref{lem:descent}, $\Pi(M)$ is
computed over $\F_4$ by a feedback system with at most $3s+O(JN)$ states.
\item By Lemma~\ref{lem:proj}(2),
$\rk_{\F_4}\Pi(B)\le\tfrac32\gamma JN=3\gamma h$.
\item By Lemma~\ref{lem:proj}(3), $\Pi(Z)$ has at most
$\eta'JN\,2^{\sqrt{\log N}/8}$ nonzero entries.
\end{itemize}
Set $\gamma=\beta/3$ and $\eta'=\eta/4$. Since $\log N=\log n+1$, we have
$2^{\sqrt{\log N}/8}\le2\cdot2^{\sqrt{\log n}/8}$. So $\Pi(Z)$ has at most
$\eta h\,2^{\sqrt{\log n}/8}$ nonzero entries. Theorem~\ref{thm:fft}(2) now
gives $3s+O(JN)\ge c\,h\sqrt{\log n}$, hence $s\ge c'JN\sqrt{\log N}$ for
all sufficiently large $N$.
\end{proof}

\begin{proof}[Proof of Corollary~\ref{cor:xor}]
For $N=2\cdot3^d$, the lower bound is Theorem~\ref{thm:main} with $B=Z=0$.
The upper bound follows from Theorem~\ref{thm:fft}(1) and the converse in
Lemma~\ref{lem:descent}. For general $N$, an XOR circuit for $H_N$ with
the last $N-2\cdot3^d$ inputs set to zero computes $H_{2\cdot3^d}$ on the
first $2\cdot3^d$ outputs, without more gates. The upper bound for $H_N$
adds $O(N)$ gates. Since $2\cdot3^d>N/3$, both bounds carry over with
modified constants.
\end{proof}

\begin{proof}[Proof of Corollary~\ref{cor:rigid}]
If $H_N=B+Z$ with $\rk B\le\gamma N$ and
$|\supp Z|\le\eta'N\,2^{\sqrt{\log N}/8}$, then Theorem~\ref{thm:main} with
$J=1$ and $M=0$ (so $s=0$) gives $0\ge c'N\sqrt{\log N}$, a contradiction.
\end{proof}

\begin{proof}[Proof of Theorem~\ref{thm:every}]
Part~1 was explained after the statement. For part~2, let
$\operatorname{char}K=2$, $n=3^d$, and let $T$ be the Fourier transform of
$(\Z/3\Z)^d$.

\emph{Case $\omega\notin K$.} Then $E=K(\omega)$ has degree two over $K$,
and it is separable since $x^2+x+1$ has derivative $1$. In the basis
$\{1,\omega\}$, multiplication by $\omega$ again acts by $W$, so the
$K$-form of $T$ is $H_N$. Lemma~\ref{lem:descent} and
Theorem~\ref{thm:fft}(1) over $E$ give both bounds.

\emph{Case $\omega\in K$.} Then $W$ has distinct eigenvalues $\omega$ and
$\omega^2$ in $K$, so $W=P\operatorname{diag}(\omega,\omega^2)P^{-1}$ for
some $P\in GL_2(K)$. Conjugating each block by $P$ and permuting
coordinates turns $H_N$ into $F\oplus\overline F$. Here $F$ is the Fourier
transform and $\overline F$ is the transform with $\omega$ replaced by
$\omega^2$, which is also a character table of $(\Z/3\Z)^d$. These changes
of basis cost $O(N)$ gates, so Theorem~\ref{thm:fft}(1) with $J=2$ gives
both bounds.
\end{proof}

\section{Discussion}
\paragraph{Circuits with AND gates.}
Our lower bounds are for XOR circuits. They do not exclude smaller
circuits for $H_N$ that also use AND gates. Whether nonlinear gates can
help compute linear maps is a well-known open problem. For bilinear maps,
the analogous descent to Boolean circuits fails. On bits $x^2=x$, so an
AND gate can create bilinear terms whose rank is not controlled by the
number of gates.

\paragraph{A consistency check.}
Alman and Williams bound the rigidity of the Walsh--Hadamard matrix from
above \cite{AW17}. Extrapolated heuristically to linear rank, their bound
allows about $N2^{O(\sqrt{\log N\log\log N})}$ changes, which is consistent
with lower bounds of the form $N2^{\Omega(\sqrt{\log N})}$.

\paragraph{Open problems.}
\begin{itemize}
\item Is $\Cx(H_N)=\Theta(N\sqrt{\log N})$? The upper bound in
Theorem~\ref{thm:fft} leaves the case $\varepsilon=0$ open.
\item The descent costs $kn$ gates for an extension of degree $k$, so it
only handles extensions of bounded degree. Do cyclic or Toeplitz maps over
$\F_2$, which arise from transforms over $\F_{2^k}$ with $k$ growing,
inherit superlinear bounds?
\item Is there an explicit matrix $A$ over $\F_2$ with
$R_A(\gamma N)\ge N^{1+\delta}$? This would meet Valiant's threshold.
\end{itemize}

\paragraph{Acknowledgments.}
The Fourier exponent theorem and its formalization are due to the
authors of \cite{FFT,Lean}.

\begin{thebibliography}{99}
\bibitem{AW17}
J.~Alman and R.~Williams. Probabilistic rank and matrix rigidity.
In \emph{Proceedings of STOC 2017}, pages 641--652, 2017.
\href{https://doi.org/10.1145/3055399.3055484}{doi:10.1145/3055399.3055484}.

\bibitem{Fri93}
J.~Friedman. A note on matrix rigidity. \emph{Combinatorica},
13(2):235--239, 1993.
\href{https://doi.org/10.1007/BF01303207}{doi:10.1007/BF01303207}.

\bibitem{Geom}
OpenAI. Euclidean geometry in the asymptotic spectrum of tensors.
Preprint, November 24, 2026.
\href{https://github.com/openai/math/blob/3344da59991ed94f9ba711217b29c1d760badfa6/preprints/Euclidean-geometry-in-the-asymptotic-spectrum-of-tensors-November-24-2026/paper.pdf}{Manuscript}.

\bibitem{FFT}
OpenAI. The logarithmic exponent of finite Fourier transforms.
Preprint, November 24, 2026.
\href{https://github.com/openai/math/blob/3344da59991ed94f9ba711217b29c1d760badfa6/preprints/The-logarithmic-exponent-of-finite-Fourier-transforms-November-24-2026/paper.pdf}{Manuscript}.

\bibitem{Lean}
OpenAI. Lean formalization: \texttt{lean/ComparatorChallenges/ExactFourier.lean},
\texttt{ExactFourier.json}, and scope document \texttt{lean/docs/379.md}.
November 24, 2026.
\href{https://github.com/openai/math/tree/3344da59991ed94f9ba711217b29c1d760badfa6/lean}{Repository}.

\bibitem{Lup}
O.~B. Lupanov. On rectifier and switching-and-rectifier schemes.
\emph{Doklady Akademii Nauk SSSR}, 111:1171--1174, 1956.

\bibitem{SSS97}
M.~A. Shokrollahi, D.~A. Spielman, and V.~Stemann. A remark on matrix
rigidity. \emph{Information Processing Letters}, 64(6):283--285, 1997.
\href{https://doi.org/10.1016/S0020-0190(97)00190-7}{doi:10.1016/S0020-0190(97)00190-7}.

\bibitem{Val77}
L.~G. Valiant. Graph-theoretic arguments in low-level complexity.
In \emph{Mathematical Foundations of Computer Science 1977}, LNCS 53,
pages 162--176, 1977.
\href{https://doi.org/10.1007/3-540-08353-7_135}{doi:10.1007/3-540-08353-7\_135}.
\end{thebibliography}
\end{document}
